\chapter{Requirements and Analysis}
\label{ch:requirements}

\section{Stakeholders and User Roles}
\label{sec:stakeholders}

The system is designed to serve three distinct actors, each with a clearly defined role in the conduct of a reading experiment.

\subsection{Researcher}
\label{sec:stakeholder-researcher}

The researcher is the primary human operator and the stakeholder around whose workflow the platform is principally organised.
In the context of a controlled reading study, the researcher is responsible for the complete experimental lifecycle: preparing reading materials and experiment templates, enrolling participants and recording their demographic attributes, ensuring the sensing hardware is correctly connected and calibrated, and conducting the session from start to finish.

During an active session, the researcher operates from a physically separate screen, observing the participant's reading behaviour through a mirrored live view without being present at the participant's display.
From this interface, the researcher must be able to monitor session health, manually apply typographic interventions, approve or reject proposals submitted by an automated decision strategy, pause and resume automated decision-making, and advance the participant through a multi-item reading sequence.
Once a session concludes, the researcher must be able to export session data for external analysis and replay recorded sessions to review gaze behaviour and intervention history at a later point.

Because the researcher coordinates every aspect of an experiment, minimising operator cognitive load is a primary concern.
A guided, step-by-step setup workflow with explicit readiness gating and a unified live control panel are therefore core requirements, not peripheral usability enhancements.

\subsection{Test Subject}
\label{sec:stakeholder-participant}

The test subject is the person recruited by the researcher to take part in a reading experiment.
Their role within the system is deliberately narrow: they read Markdown-formatted text presented in the reading interface while their gaze behaviour is recorded.
The test subject does not interact with system configuration, session management, or any researcher-facing controls.

During a session, the test subject may be exposed to typographic adaptations --- changes to font size, line width, line height, letter spacing, font family, colour palette, or theme mode --- applied by the researcher or by an automated strategy.
A critical requirement arising from this role is that such adaptations must not disrupt reading flow or cause the test subject to lose their place in the text, as perceptible disruption would confound the experimental results.
Following each reading item, the test subject may also be asked to answer comprehension quiz questions, the results of which are recorded as part of the session data.

\subsection{External Module Provider}
\label{sec:stakeholder-provider}

The external module provider is an out-of-process software agent that connects to the platform, declares its capabilities, receives structured context containing gaze observations and session state, and returns commands such as intervention recommendations or derived fixation events.

A central research motivation for this platform is the ability to evaluate different decision and analysis strategies under comparable experimental conditions.
This actor motivates the need for a stable and documented integration protocol so that different strategies --- rule-based, machine-learning-based, or otherwise --- can be substituted without changes to the core system.
The platform must support a range of execution modes, from fully autonomous application of proposals to researcher-gated approval, so that the degree of automation can itself be varied as an experimental condition.
The system must also continue operating normally if a provider disconnects mid-session, ensuring that a provider failure does not abort an ongoing experiment.

\section{Use Cases and Scenarios}
% Primary scenarios: run a reading study, replay a session, swap an
% intervention module, log gaze data, etc.

\section{Functional Requirements}
% FR-1, FR-2, ... in a table with priority (MoSCoW or similar).

\section{Non-Functional Requirements}
% Performance (real-time latency budgets), reliability, modularity,
% experimentability, observability, privacy, portability.

\section{Constraints}
% Hardware (eye-tracker capabilities), monorepo / deployment, ethics.

\section{Summary}
% Brief recap, leading into the architecture chapter.
